\documentclass[titlepage]{article}

\usepackage{titlepic}
\usepackage{tikz}

\title{Groundwater Recharge in Semi-Arid Basins}
\author{N. Farouk \\ Department of Hydrology}
\date{\today}

% titlepic normally holds a picture; here it holds a drawn vertical
% rule flanked by short horizontal ticks, a graphic device rather
% than a photograph, giving the title page a distinct look from the
% other title-page starters in the catalogue.
\titlepic{%
\begin{tikzpicture}
  \draw[line width=0.6pt] (0,0) -- (0,3.2);
  \draw[line width=0.6pt] (-0.15,0) -- (0.15,0);
  \draw[line width=0.6pt] (-0.15,3.2) -- (0.15,3.2);
\end{tikzpicture}%
}

\begin{document}

\maketitle

\begin{abstract}
Groundwater recharge in semi-arid basins depends strongly on the
timing of rainfall relative to soil saturation, not merely on total
annual rainfall. This report presents recharge estimates for the
Kavango Basin derived from a network of 14 monitoring boreholes over
a three-year period, and finds that recharge is dominated by a small
number of high-intensity storm events rather than by cumulative
seasonal rainfall.
\end{abstract}

\section{Introduction}

Conventional recharge estimates in semi-arid regions often assume a
roughly linear relationship between seasonal rainfall total and
aquifer recharge. Borehole data from the Kavango Basin, collected
over three wet seasons (2023--2026), suggests this assumption breaks
down once rainfall is disaggregated by storm intensity rather than
treated as a seasonal total.

\section{Method}

Water-table depth was logged hourly at 14 boreholes distributed
across the basin's three principal soil types (sandy loam, clay-rich
floodplain, and shallow laterite). Recharge events were identified as
water-table rises exceeding the borehole's measurement noise floor
within a 48-hour window following rainfall, and attributed to the
specific storm event whose timing matched the rise.

\end{document}
